% Chapter 4: The Four Properties
% DEFINES: ch:four-properties; sec:totality-property, sec:uniformity,
%   sec:correctness, sec:composability, sec:parameter-tuple, sec:ch4-notes;
%   def:totality, def:uniformity, def:correctness, def:composability,
%   thm:composition, ex:delta-watchlist, ex:composition, rem:deniability,
%   tab:compounding, tab:parameters
% REFERENCES (backward): sec:untrusted-sees, sec:measured-not-negligible (Part I);
%   def:cipher-map, sec:totality-support, ex:watchlist-numbers (Ch 3)
% REFERENCES (forward, part labels): part:confidentiality, part:algebra,
%   part:constructions

\chapter{The Four Properties}
\label{ch:four-properties}

\Cref{ch:cipher-maps} gave the cipher map as a tuple and named, in passing,
the structural facts that make it useful: it is total, its tokens can be
made to look alike, it is correct up to an error budget, and it composes.
This chapter states those four as numbered definitions and attaches to each
the parameter that measures it. Together they are the contract a
construction must meet, and they are the book's definitions of record: the
later parts measure, build, and compose against exactly these.

The four properties also pay off a promise. \Cref{sec:untrusted-sees}
listed four things the untrusted machine cannot do, decode a token,
distinguish a real query from filler, recover the domain, or be sure a
result is correct, and tied each to a property still to be defined. Each
definition below supplies one of them, and the connection is drawn as it
arrives. The parameters are introduced in the spine's canonical order
$(\eta, \varepsilon, \mu, \delta)$, and where this chapter and the standalone
papers differ on a definition, the chapter follows the reconciled spine and
the divergence is logged for back-port.

\section{Totality}
\label{sec:totality-property}

The first property is the one \Cref{sec:totality-support} already leaned on:
$\fhat$ answers every string.

\begin{definition}[Totality, Property 1]
\label{def:totality}
A cipher map $(\fhat, \enc, \dec, s)$ is \emph{total}: $\fhat$ is defined on
all of $\B^n$, and for a token $c \notin \mathrm{im}(\enc)$, that is, a
string encoding no input, the output $\fhat(c)$ is a uniform random element
of $\B^n$ under the random oracle model. Equivalently, every $c \in \B^n$
yields $\fhat(c) \in \B^n$, with no distinguished ``undefined'' result.
\end{definition}

Two consequences follow, one for each party. For the trusted machine,
totality is what makes the noise-decode probability $\varepsilon$ of
\Cref{sec:totality-support} a usable knob: a uniformly random token decodes
to a value rather than to $\bot$ with probability $\varepsilon$, and the
construction sets $\varepsilon$ by how much of the token space it treats as
meaningful. For the untrusted machine, totality is the first privacy
guarantee. Because a real query $\enc(x, k)$ and a random filler string both
produce $n$-bit answers, and nothing distinguishes the two at evaluation
time, the untrusted machine cannot tell which strings in a query stream
carry meaning. This is the first of the four cannots of
\Cref{sec:untrusted-sees}, now a structural feature of $\fhat$ rather than a
hope: in the watchlist of \Cref{ex:watchlist-numbers}, the trusted machine
may pad its lookups with uniformly random $12$-bit strings, and since each
returns an answer, the server cannot subtract the filler back out.

\section{Representation Uniformity}
\label{sec:uniformity}

The second property is the subtlest, and it is the one the rest of the book
leans on hardest, so this section states it with care. It is also the
chapter's first correction to the standalone formalism: the definition is
\emph{image-relative}, measured against the populated support, not the
ambient token space.

Recall from \Cref{sec:tuple} that a value $x$ may be encoded as any of
$K(x) \geq 1$ tokens. Fix a query distribution $D$ on $X$, the demand. A
query draws $x \sim D$ and then a representation index $k$ uniformly from
$x$'s $K(x)$ tokens, inducing a distribution on tokens
\[
  Q(c) \;=\; \sum_{x \in X} D(x) \cdot
             \frac{\lvert \{k : \enc(x, k) = c\} \rvert}{K(x)},
\]
so that a token $c$ encoding $x$ carries probability $D(x)/K(x)$, the demand
for $x$ spread across its representations. The observed token traffic is a
sample from $Q$, supported on $\mathrm{im}(\enc)$.

\begin{definition}[Representation uniformity, Property 2]
\label{def:uniformity}
A cipher map has \emph{representation uniformity $\delta$} if the induced
token distribution $Q$ is $\delta$-close, in total variation, to the uniform
distribution on the populated support:
\[
  \delta \;=\; \TV\!\bigl(Q,\; U_{\mathrm{im}}\bigr),
  \qquad
  U_{\mathrm{im}} = \mathrm{Uniform}\bigl(\mathrm{im}(\enc)\bigr),
  \qquad
  H^{*} = \log_2 \lvert \mathrm{im}(\enc) \rvert .
\]
The reference is the uniform distribution on $\mathrm{im}(\enc)$, the tokens
that actually occur, and $H^{*}$ is the most entropy the token traffic can
carry. The reference is \emph{not} the uniform distribution on the ambient
space $\B^n$.
\end{definition}

The image-relative reference is the whole point of the definition, so it is
worth saying why the ambient alternative fails. A cipher map's tokens
occupy a vanishing fraction of $\B^n$, ten of $4096$ strings in
\Cref{ex:watchlist-numbers}. Measured against uniform on all of $\B^n$, even
a perfectly flattened $Q$ sits at total variation almost $1$ from the
reference simply because it is supported on a sparse set, so an ambient
$\delta$ would be near its maximum for every realizable construction and
would distinguish nothing.\footnote{This is why \Cref{sec:totality-support}
insisted on the distinction between the populated support and the ambient
space: confidentiality is closeness to uniform \emph{on the support}.
Measured ambiently the bound is unsatisfiable for any sparse image.}
Confidentiality must be measured against the support the traffic actually
lives on, and $\delta$ as defined is exactly that distance.

\paragraph{The construction principle.}
Small $\delta$ is bought by \emph{multiple representations}. Assign
$K(x) \propto D(x)$, so that frequently queried values get \emph{more}
tokens, classical homophonic substitution, and the per-token probability
$D(x)/K(x)$ flattens toward a constant. The direction matters and is a place
the historical record went wrong: more demand earns more representations, so
that no single token is queried often. The inverted prescription
$K(x) \propto 1/D(x)$, which would give the most-queried value the
\emph{fewest} tokens and concentrate rather than spread its traffic, is
wrong and was corrected in the formalism. With the right allocation the
homophonic-allocation identity makes the bound exact:
\[
  \delta \;=\; \TV(Q, U_{\mathrm{im}}) \;=\; \TV(D, K/N),
  \qquad N = \sum_x K(x),
\]
where $K/N$ is the fraction of the token budget each value receives; the
right-hand side is bounded by the tight $(\lvert X \rvert - 1)/N$. When
$K \propto D$ exactly, $K/N = D$ and $\delta = 0$.

\begin{example}[Flattening the watchlist]
\label{ex:delta-watchlist}
Take the four names with demand $D = (0.4, 0.2, 0.2, 0.2)$, alice the
frequent one. With a single representation each ($K = (1,1,1,1)$, $N = 4$),
the token distribution is $Q = D$ and the support is four tokens, so
\[
  \delta = \TV\bigl((0.4,0.2,0.2,0.2),\,(\tfrac14,\tfrac14,\tfrac14,\tfrac14)\bigr)
         = \tfrac12\bigl(0.15 + 3 \cdot 0.05\bigr) = 0.15 .
\]
Now allocate homophonically, $K = (4, 2, 2, 2)$ as in
\Cref{ex:watchlist-numbers}, so $N = 10$ and $K/N = (0.4, 0.2, 0.2, 0.2) =
D$. Each of alice's four tokens now carries $0.4/4 = 0.1$, each other token
$0.2/2 = 0.1$: the ten-token distribution is exactly uniform and
$\delta = \TV(D, K/N) = 0$. The most-queried value took the most tokens,
the opposite of the inverted rule, and the traffic flattened completely.
\end{example}

This is the parameter the measurable-not-negligible thread of
\Cref{sec:measured-not-negligible} becomes. Small $\delta$ is what lets the
untrusted machine learn little about \emph{which} value a token encodes, and
it is the lever the confidentiality theory turns into a quantitative leakage
bound.\footnote{The bridge from $\delta$ to a confidentiality measure, the
entropy ratio and its Fannes-type lower bound, is the subject of
\Cref{part:confidentiality}.} It pays off the third of the four cannots: with
$Q$ near uniform on the support, the token frequencies betray nothing about
the domain. One honest limit, carried forward: $\delta$ bounds the
\emph{marginal} distribution of tokens. It says nothing about joint
distributions, so correlations between separately encoded values can still
leak, a matter the encoding-granularity discussion of \Cref{part:algebra}
takes up.

\section{Correctness}
\label{sec:correctness}

The third property is the error budget, and it is where the approximation in
``trapdoor approximation'' is quantified.

\begin{definition}[Correctness, Property 3]
\label{def:correctness}
A cipher map is \emph{$\eta$-correct} if it computes the latent function
correctly on in-domain inputs with probability at least $1 - \eta$: for
$x$ drawn uniformly from $X$ and $k$ uniformly from $x$'s representations,
\[
  \Pr_{x, k}\bigl[\dec\bigl(\fhat(\enc(x, k))\bigr) \neq f(x)\bigr]
  \;\leq\; \eta .
\]
\end{definition}

The parameter $\eta$ is a construction cost, set by how many constraints the
secret $s$ must satisfy: tolerating a larger $\eta$ makes the seed search
faster and the resulting structure smaller, a trade made precise in
\Cref{part:constructions}. Two features matter for the model. First, once a
cipher map is built, its failures are deterministic; the same in-domain
elements fail for a given $s$, so $\eta$ is a frequency over $X$, not a
fresh coin on each call. Second, $\eta$ is a guarantee to the trusted
machine only. The untrusted machine never decodes, so it never sees whether
a result was right; correctness is the trusted party's concern, not a signal
the server can read.\footnote{The error rate $\eta$ is the false-negative
rate of the Bernoulli error model that the hash-based construction inherits;
the model, and why per-element errors compose tractably, are summarized in
\Cref{app:bernoulli}.}

\begin{remark}[Approximation as privacy]
\label{rem:deniability}
A nonzero error rate is not only a cost; it also buys privacy, and this pays
off the last of the four cannots of \Cref{sec:untrusted-sees}. When
$\eta > 0$, an answer carries ambiguity: a reported membership might be a
true hit or a cipher-map error, and nothing in the result distinguishes the
two. With $\eta = 0$ the answer is exact and an adversary who learns it
learns the truth; with $\eta > 0$ each answer is deniable. This is the dual
of representation uniformity. Property 2 hides \emph{which} value was
queried; a nonzero $\eta$ hides \emph{what the answer means}.
\end{remark}

\section{Composability and the Composition Theorem}
\label{sec:composability}

The fourth property is what makes cipher maps a calculus rather than a
collection of isolated lookups: they chain, and the error of a chain is
predictable from the errors of its links.

\begin{definition}[Composability, Property 4]
\label{def:composability}
Cipher maps compose. If $\fhat$ is a cipher map for $f : X \to Y$ and $\ghat$
is a cipher map for $g : Y \to Z$, sharing the token space $\B^n$, then
$\ghat \circ \fhat$, defined by $(\ghat \circ \fhat)(c) = \ghat(\fhat(c))$,
is a cipher map for $g \circ f$. Totality is what permits this: the output
token of $\fhat$ is a well-formed input to $\ghat$, whether or not it is a
codeword $\ghat$ recognizes.
\end{definition}

The reason composition is useful is that its error is not merely bounded but
\emph{computable} from the parts.

\begin{theorem}[Composition]
\label{thm:composition}
Let $\fhat_1, \ldots, \fhat_m$ be cipher maps with correctness parameters
$\eta_1, \ldots, \eta_m$, chained so that each receives the previous one's
output. The chain $\fhat_m \circ \cdots \circ \fhat_1$ is a cipher map for
the composite latent function with correctness parameter
\[
  \eta_{\mathrm{total}} \;\leq\; 1 - \prod_{i=1}^{m} (1 - \eta_i) .
\]
\end{theorem}

\begin{proof}
Let $C_i$ be the event that stage $i$ produces the correct output given a
correct input. The chain produces the correct final output whenever every
stage does: if each stage receives a correct intermediate value and answers
correctly, the composite is correct. So $\bigcap_{i} C_i$ implies the chain
is correct, and therefore
\[
  \Pr[\text{chain correct}] \;\geq\; \Pr\Bigl[\bigcap_{i} C_i\Bigr]
  \;=\; \prod_{i=1}^{m} \Pr\bigl[C_i \mid C_1, \ldots, C_{i-1}\bigr]
  \;\geq\; \prod_{i=1}^{m} (1 - \eta_i),
\]
the equality by the chain rule and the last step because, conditioned on
every upstream stage being correct, stage $i$ receives its intended input
and so errs with probability at most $\eta_i$. Taking complements,
$\eta_{\mathrm{total}} = \Pr[\text{chain incorrect}] \leq 1 - \prod_i (1 -
\eta_i)$.
\end{proof}

The bound is an \emph{inequality}, and the gap is real, not slack in the
argument. ``Every stage correct'' is sufficient for a correct chain but not
necessary: a stage can err and a later stage can map the corrupted token
back to a correct value, masking the error. The cipher map is a hash, and a
valid-looking intermediate token does not certify a valid input, so
cancellations occur and the true error can fall strictly below
$1 - \prod_i(1 - \eta_i)$. Stating the relation as an equality, as a casual
reading of the two-stage formula
$\eta_{g \circ f} = \eta_f + \eta_g - \eta_f \eta_g$ invites, overcounts
exactly these masked errors; the relation is $\leq$.

For small per-stage rates the bound linearizes. Expanding the product,
$1 - \prod_i(1 - \eta_i) \approx \sum_i \eta_i$ to first order, so a short
chain of accurate maps accumulates error roughly additively, the same
content as the union bound. The predictability is the message of the
composition-predictability thread: a designer who knows the per-stage
budgets knows the chain's budget, with no construction-specific surprises.

\begin{example}[A three-stage chain]
\label{ex:composition}
Compose three cipher maps with $\eta = (0.01, 0.02, 0.01)$. The theorem
gives
\[
  \eta_{\mathrm{total}} \leq 1 - (0.99)(0.98)(0.99) = 1 - 0.9605 = 0.0395,
\]
slightly below the naive sum $0.04$, the difference being the
doubly-counted joint failures the product bound removes. The chain is at
least $96\%$ correct, and a designer reads that off the three budgets
directly.
\end{example}

Errors do compound as chains deepen, but slowly and predictably.
\Cref{tab:compounding} traces a chain of identical $99\%$-correct maps:
the chain stays above $90\%$ correct out to ten stages and is still right
roughly a third of the time at a hundred. The growth is sub-linear in $m$
until $m\eta$ approaches one, which is what makes deep cipher-map programs
feasible at all.

\begin{table}[t]
\centering
\caption{Compounded error $\eta_{\mathrm{total}} = 1 - (1-\eta)^m$ for a
chain of $m$ identical maps with per-stage $\eta = 0.01$.}
\label{tab:compounding}
\begin{tabular}{@{}rcccccc@{}}
\toprule
$m$ & 1 & 2 & 5 & 10 & 50 & 100 \\
\midrule
$\eta_{\mathrm{total}}$ & 0.010 & 0.020 & 0.049 & 0.096 & 0.395 & 0.634 \\
\bottomrule
\end{tabular}
\end{table}

\section{The Parameter Tuple}
\label{sec:parameter-tuple}

The four properties contribute four parameters, and collecting them gives
the tuple $(\eta, \varepsilon, \mu, \delta)$ that characterizes a cipher map.
Three have appeared; the fourth, $\mu$, is the cost of the values
themselves. \Cref{tab:parameters} sets them side by side in the spine's
canonical order.

\begin{table}[t]
\centering
\caption{The parameter tuple, in canonical order.}
\label{tab:parameters}
\begin{tabular}{@{}clll@{}}
\toprule
Symbol & Name & Controls & Set by \\
\midrule
$\eta$ & Correctness & error rate on in-domain inputs & seed search / structure \\
$\varepsilon$ & Noise decode & chance a random token decodes to a value & meaningful fraction of $\B^n$ \\
$\mu$ & Value cost & bits to encode an output value, $H(Y)$ & the codomain $Y$ \\
$\delta$ & Uniformity & TV of $Q$ from uniform on $\mathrm{im}(\enc)$ & representation multiplicity $K(x)$ \\
\bottomrule
\end{tabular}
\end{table}

The value cost $\mu = H(Y)$ is the entropy of the output distribution, the
irreducible number of bits an answer must carry; a Boolean watchlist has
$\mu \leq 1$, a richer codomain more. It joins $-\log_2 \varepsilon$, the
cost of hiding values among noise, to set a cipher map's space, a duality
that \Cref{part:constructions} derives. The four parameters are largely
independent: a construction can hit one without the others, and the next
chapter's HashSet, for instance, is perfectly correct ($\eta = 0$) yet has
no representation uniformity at all ($K(x) = 1$, so $\delta$ does not apply).
The properties are a contract with four separable terms, not a single dial.

\section{Notes and Provenance}
\label{sec:ch4-notes}

The four properties and the tuple $(\eta, \varepsilon, \mu, \delta)$ are the
spine's, in its order. Two definitions here correct the standalone
formalism and are flagged for back-port. First, representation uniformity
(\Cref{def:uniformity}) is stated \emph{image-relative}, against uniform on
$\mathrm{im}(\enc)$ with $H^{*} = \log_2\lvert\mathrm{im}(\enc)\rvert$; the
spine's Property 2 headline still measures against the ambient $\B^n$, which
is unsatisfiable for sparse images, and the corrected form is the one the
confidentiality paper and this book use. Second, the composition theorem
(\Cref{thm:composition}) is stated and proved as an \emph{inequality}; an
error in one stage can be masked by another, so the equality form that the
spine's headline and an earlier paper draft carried overcounts. The proof is
elementary, and the same bound is reached independently from the Bernoulli
composition theorem, a noisy-gates analysis, and a probabilistic union
bound, which is some assurance it is structural rather than an artifact of
one derivation. The homophonic construction principle $K(x) \propto D(x)$,
and the warning against its inverted form, are carried from the corrected
spine; the deniability reading of nonzero $\eta$ (\Cref{rem:deniability}) is
the spine's. The Bernoulli error model behind $\eta$ is summarized in
\Cref{app:bernoulli}.
